\documentclass[tikz,border=8pt]{standalone}
\usepackage{tikz}
\usepackage{amsmath}
\usepackage{ctex}
\begin{document}
\begin{tikzpicture}
  % Discrete CDF example: P(1)=0.2, P(2)=0.3, P(3)=0.3, P(4)=0.2

  % Axes
  \draw[->] (-0.5,0) -- (5.5,0) node[right,font=\small] {$x$};
  \draw[->] (0,-0.3) -- (0,5.2) node[above,font=\small] {$F(x)$};

  % y-axis ticks
  \foreach \y/\lbl in {0/0, 1.0/0.2, 2.5/0.5, 4.0/0.8, 5.0/1.0} {
    \draw (-0.1,\y) -- (0,\y);
    \node[font=\scriptsize,left] at (-0.12,\y) {\lbl};
  }

  % x-axis labels
  \foreach \x in {1,2,3,4} {
    \node[font=\small,below] at (\x,-0.12) {\x};
    \draw (\x,-0.06) -- (\x,0.06);
  }

  % Staircase CDF segments
  % F = 0 for x<1
  \draw[blue,ultra thick] (-0.5,0) -- (1,0);
  % F = 0.2 for 1<=x<2
  \draw[blue,ultra thick] (1,1.0) -- (2,1.0);
  % F = 0.5 for 2<=x<3
  \draw[blue,ultra thick] (2,2.5) -- (3,2.5);
  % F = 0.8 for 3<=x<4
  \draw[blue,ultra thick] (3,4.0) -- (4,4.0);
  % F = 1.0 for x>=4
  \draw[blue,ultra thick] (4,5.0) -- (5.5,5.0);

  % Jump points (closed circles at right)
  \foreach \x/\y in {1/1.0, 2/2.5, 3/4.0, 4/5.0} {
    \fill[blue] (\x,\y) circle (3pt);
  }
  % Open circles at left of each jump
  \foreach \x/\y in {1/0, 2/1.0, 3/2.5, 4/4.0} {
    \draw[blue,fill=white] (\x,\y) circle (3pt);
  }

  % Dashed lines for jump heights
  \foreach \x/\y in {1/1.0, 2/2.5, 3/4.0, 4/5.0} {
    \draw[gray,dashed] (\x,0) -- (\x,\y);
    \draw[gray,dashed] (0,\y) -- (\x,\y);
  }

  % Jump size annotations
  \draw[<->,red,thin] (1.05,0) -- (1.05,1.0) node[midway,right,font=\scriptsize,red] {$+0.2$};
  \draw[<->,red,thin] (2.05,1.0) -- (2.05,2.5) node[midway,right,font=\scriptsize,red] {$+0.3$};
  \draw[<->,red,thin] (3.05,2.5) -- (3.05,4.0) node[midway,right,font=\scriptsize,red] {$+0.3$};
  \draw[<->,red,thin] (4.05,4.0) -- (4.05,5.0) node[midway,right,font=\scriptsize,red] {$+0.2$};

  % Title
  \node[font=\small\bfseries] at (2.5,5.6) {离散型 CDF（阶梯函数）};

  % Annotation: right-continuous
  \node[font=\scriptsize,fill=yellow!20,draw,rounded corners=2pt,align=left,
        anchor=south east]
    at (5.4,0.6) {右连续：\\$F(x_k)$取右极限\\空心 = 左极限};

\end{tikzpicture}
\end{document}
